\chapter{Objetivos y metodología de trabajo}
\label{cap:4}

\section{Introducción al capítulo}

Este capítulo conecta el problema descrito en los capítulos anteriores con la estrategia seguida para resolverlo. El TFM no plantea una herramienta generica de emergencias, sino un prototipo software de apoyo a la decisión para priorización temprana de incidentes 112 en Castilla y León. La línea de trabajo se articula alrededor de datos reales abiertos, una etiqueta académica P1--P4 construida mediante supervisión débil, una Capa 2 interpretable basada en RuleFit-lite y una Capa 3 orientada a explicar la recomendación con trazabilidad normativa.

La finalidad del capítulo es fijar con claridad qué se pretende conseguir, qué preguntas guían el desarrollo, qué aportación de inteligencia artificial realiza el trabajo y cómo se ha organizado el equipo para avanzar de forma incremental sin perder coherencia entre prototipo, evaluación y memoria.

\section{Objetivo general}

El objetivo general de este Trabajo Fin de Máster es:

\textit{Diseñar, implementar y evaluar un prototipo software modular de apoyo a la decisión para la priorización inicial de incidentes 112 en Castilla y León, capaz de transformar texto operativo y datos parciales en una recomendación P1--P4 explicable, trazable y revisable, mediante una arquitectura de tres capas que combine extracción NLP, aprendizaje interpretable con RuleFit-lite, explicación normativa local y supervisión humana obligatoria.}

El objetivo es específico porque delimita dominio, territorio, usuario y salida esperada. Es medible porque la Capa 2 se evalúa mediante métricas reproducibles como macro-F1, sensibilidad P1, matriz de confusión, falsos negativos críticos, estabilidad temporal y análisis por subgrupos. Es alcanzable porque se apoya en datos abiertos del Registro 112 Castilla y León, procedimientos reproducibles y un prototipo académico acotado. Es relevante porque aborda una decisión de alto impacto en emergencias sin sustituir al operador. Y es temporalmente delimitado porque se concreta en una evaluación cerrada dentro del alcance del TFM.

\section{Objetivos específicos}

El objetivo general se descompone en ocho objetivos específicos:

\begin{enumerate}
    \item[OE1.] Analizar el problema de la priorización temprana de incidentes 112, identificando necesidades del operador, restricciones de información y riesgos de automatización en un dominio crítico.
    \item[OE2.] Revisar el marco normativo estatal y autonómico aplicable a Castilla y León, extrayendo criterios de gravedad, coordinación, escalado y trazabilidad que fundamentan el sistema.
    \item[OE3.] Definir la taxonomia operativa, la escala académica P1--P4 y las variables V01--V15, separando variables disponibles en primera decisión de variables contextuales diferidas o potencialmente sujetas a \textit{leakage}.
    \item[OE4.] Auditar el Registro de Emergencias 112 de Castilla y León, construir el conjunto de datos limpio, documentar limitaciones y generar etiquetas académicas P1--P4 mediante supervisión débil.
    \item[OE5.] Diseñar e implementar una Capa 2 interpretable, comparar una línea base heurística con RuleFit-lite y una ejecución diagnóstica de \texttt{imodels}, seleccionar el modelo mediante la partición de validación y reservar la partición de prueba para la estimación final.
    \item[OE6.] Diseñar una Capa 3 de explicación que convierta reglas, probabilidades, evidencias y citas normativas en una justificación comprensible para el operador, sin delegar la decisión en un LLM.
    \item[OE7.] Integrar las capas mediante contratos versionados, servicio web auditable, respuesta alternativa determinista, registros de inferencia y mecanismos de retroalimentación humana.
    \item[OE8.] Evaluar el prototipo con una partición estratificada y otra temporal, control de fugas de información, estudio de errores P1, análisis por subgrupos, revisión interna de una muestra de casos, evaluación automatizada de fidelidad explicativa y evaluación exploratoria con participantes.
\end{enumerate}

En relación con OE5, la Capa 2 queda preparada para su consumo por Capa 3 mediante una salida explicable que incluye prioridad recomendada, reglas activadas, detalles de decisión y contexto estructurado para explicación. Esta salida permite que Capa 3 genere una explicación operativa sin recalcular ni sustituir la prioridad asignada por el modelo interpretable.

En relación con OE8, los tres autores han revisado conjuntamente una muestra de 119 casos, compuesta por 59 falsos negativos P1 y 60 casos equilibrados adicionales. Esta revisión interna es cualitativa y no equivale a tres anotaciones independientes ni a una validación externa con operadores del 112. De forma separada, la fidelidad de las explicaciones se evalúa automáticamente sobre esos mismos casos mediante un juez determinista reproducible. La evaluación se completa con un estudio exploratorio de nueve participantes, cuyos resultados describen utilidad percibida y comprensibilidad sin pretensión de representatividad estadística.

Esta secuencia mantiene una progresión lógica: primero se delimita el dominio, después se construye una base empírica defendible, a continuación se implementa el modelo interpretable y finalmente se evalúa su comportamiento desde una perspectiva técnica, explicativa y humana.

\begin{table}[htbp]
\centering
\scriptsize
\setlength{\tabcolsep}{3pt}
\begin{tabular}{|p{1.1cm}|p{4.4cm}|p{4.9cm}|p{3.2cm}|}
\hline
\textbf{Obj.} & \textbf{Evidencia principal} & \textbf{Soporte documental} & \textbf{Criterio de cumplimiento} \\ \hline
OE1 & Planteamiento del problema, usuario y riesgos de automatización & Capítulos 1 y 4 & Dominio, alcance y supervisión humana delimitados \\ \hline
OE2 & Marco estatal, autonómico y europeo & Capítulo 3; Tabla de fuentes normativas & Normas vinculadas a variables, reglas y límites \\ \hline
OE3 & Taxonomía, variables V01--V15 y política anti-\textit{leakage} & Anexos A y B; Capítulo 7 & Variables disponibles y prohibidas documentadas \\ \hline
OE4 & Conjunto limpio, auditoría, particiones y etiquetas débiles & Capítulo 7; Anexo C & Conjunto reproducible y etiqueta académica trazable \\ \hline
OE5 & Línea base heurística, RuleFit-lite y comparativa con alternativa externa & Capítulo 9; Anexo E & Modelo seleccionado en validación y salida consumible por Capa 3 \\ \hline
OE6 & Explicación con LLM local, RAG, MCP y respuesta alternativa determinista & Capítulos 6 y 8; Anexo L & Capa 3 explica sin recalcular prioridad \\ \hline
OE7 & Integración Capa 1--Capa 2--Capa 3 y prueba de funcionamiento & Capítulos 8 y 9 & 5/5 escenarios integrados y 4/4 casos con LLM local \\ \hline
OE8 & Evaluación final, errores P1, revisión interna, fidelidad automatizada y estudio exploratorio & Capítulo 9; Anexo D & 59 falsos negativos P1, 119 casos revisados y nueve participantes \\ \hline
\end{tabular}
\caption{Relación entre objetivos específicos, evidencias y criterios de cumplimiento.}
\label{tab:objetivos-evidencias-cumplimiento}
\end{table}

\begin{center}
\textit{Fuente: elaboración propia.}
\end{center}

\section{Preguntas de investigación}

Las preguntas de investigación orientan el diseño y la evaluación del sistema:

\begin{enumerate}
    \item[PI1.] ¿Puede construirse una etiqueta académica P1--P4 defendible cuando el conjunto abierto no contiene prioridad oficial?
    \item[PI2.] ¿Qué aporta un modelo interpretable como RuleFit-lite frente a una línea base heurística basada en señales textuales y reglas operativas?
    \item[PI3.] ¿Es posible mantener una sensibilidad P1 suficiente sin convertir el sistema en una herramienta excesivamente conservadora que degrade el rendimiento global?
    \item[PI4.] ¿Hasta qué punto los resultados se mantienen estables en una partición temporal posterior?
    \item[PI5.] ¿Qué errores críticos aparecen, destacando falsos negativos P1, y qué evidencia debe mostrarse para que puedan ser revisados por humanos?
    \item[PI6.] ¿Las explicaciones generadas por Capa 3 son fieles a las reglas, probabilidades y evidencias producidas por Capa 2?
\end{enumerate}

Estas preguntas evitan presentar el TFM como una promesa de automatización total. El foco está en demostrar si una arquitectura responsable puede producir recomendaciones útiles, explicables y auditables dentro de un prototipo académico.

\section{Contribución en inteligencia artificial}

La contribución de IA comprende la transformación de información operativa en
variables estructuradas, la construcción de etiquetas mediante supervisión débil,
la comparación de alternativas interpretables y la separación entre priorización y
explicación. La evaluación se respalda con métricas, evidencias reproducibles, una
muestra revisada internamente y un análisis automatizado de fidelidad.

La selección de RuleFit-lite se justifica por su rendimiento en validación, su sensibilidad ante P1 y la posibilidad de exportar reglas. La implementación canónica alternativa se emplea únicamente como contraste diagnóstico con una muestra de ajuste reducida, por lo que sus costes no permiten establecer una comparación computacional en condiciones equivalentes. La partición de prueba se reserva para estimar el rendimiento final del modelo ya seleccionado.

\section{Enfoque metodológico}

El desarrollo sigue un enfoque iterativo e incremental adaptado a un TFM grupal. El
equipo ha trabajado por bloques trazables: definición del dominio, datos, arquitectura,
Capa 2, Capa 3, integración, evaluación y redacción. Este enfoque permitió
contrastar la línea base heurística con RuleFit-lite y una alternativa externa de la familia RuleFit.

La metodología adopta una lógica incremental: cada fase genera una evidencia verificable, se contrasta con el objetivo correspondiente y se integra en la secuencia experimental. Este enfoque resulta adecuado porque el TFM combina investigación aplicada, desarrollo de un prototipo y redacción académica.

La metodología prioriza la evidencia reproducible frente a afirmaciones operativas no verificadas. Cada decisión relevante se vincula con una evidencia académica: conjunto limpio, guía de etiquetado, particiones, línea base heurística, RuleFit-lite, evaluación final y pruebas integradas. Esta trazabilidad permite defender que los resultados no dependen de una narración posterior, sino de un flujo experimental documentado.

\section{Fases de desarrollo}

El trabajo se ha estructurado en siete fases:

\begin{enumerate}
    \item \textbf{Delimitación del dominio.} Análisis del problema 112, marco normativo estatal y autonómico, alcance territorial CyL y principios de supervisión humana.
    \item \textbf{Datos y variables.} Auditoría del Registro 112 CyL, definición de V01--V15, identificación de columnas prohibidas por \textit{leakage} y construcción del conjunto de datos limpio.
    \item \textbf{Supervisión débil.} Generación de etiquetas académicas P1--P4, medición de acuerdo, ablación anti-circularidad y construcción de particiones estratificada y temporal.
    \item \textbf{Capa 2 interpretable.} Implementación de la línea base heurística, RuleFit-lite y una alternativa externa; extracción de reglas, pesos, probabilidades y métricas por clase.
    \item \textbf{Capa 3 explicativa e integración.} Arquitectura explicativa, trazabilidad, registro auditable e integración del flujo de inferencia.
    \item \textbf{Evaluación y trazabilidad.} Evaluación final, análisis por subgrupos, errores P1, revisión interna conjunta, fidelidad automatizada y síntesis de evidencias.
    \item \textbf{Redacción y coherencia final.} Alineación de capítulos, anexos, métricas, limitaciones y trabajo futuro para que la defensa refleje exactamente el sistema construido.
\end{enumerate}

\section{Consideraciones eticas}

El sistema se diseña para un dominio de alto impacto, por lo que incorpora salvaguardas desde el principio. La primera es la supervisión humana: el prototipo recomienda, pero no moviliza recursos ni activa planes. La segunda es la explicabilidad: toda recomendación debe acompañarse de reglas, probabilidades, evidencias y, cuando proceda, citas normativas. La tercera es el control anti-\textit{leakage}: se excluyen campos posteriores a la decisión para no entrenar un modelo que aprenda consecuencias del operativo. La cuarta es la soberanía del dato: la inferencia se concibe en local y sin dependencia de APIs externas para datos sensibles.

También se documenta el uso de herramientas de IA generativa como apoyo a la elaboración del TFM. Dichas herramientas se han empleado para organización, redacción preliminar y revisión de coherencia, pero las decisiones técnicas, metodológicas y académicas corresponden al equipo y se contrastan con las evidencias reproducibles del proyecto.

\section{Conclusiones del capítulo}

Los objetivos, preguntas y metodología definidos en este capítulo fijan una línea de trabajo coherente con el resto de la memoria: un prototipo académico, territorialmente centrado en Castilla y León, que usa datos reales abiertos, supervisión débil, un modelo interpretable y explicaciones trazables. Esta formulación permite que los capítulos posteriores no dependan de promesas abstractas, sino de evidencias concretas: conjunto limpio, particiones, modelos, métricas, revisión preparada y evaluación final.
